\documentclass{beamer}
\usetheme{boxes}
\usecolortheme{default}
\usepackage[T1]{fontenc}
\usepackage{amsmath,amssymb}
\usepackage{csquotes}
\usepackage{array}
\usepackage{colortbl}
\setbeamertemplate{navigation symbols}{}


% CUSTOM FOOTER CONFIGURATION (Page number / Total pages)
\setbeamertemplate{footline}{%
  \leavevmode%
  \hbox{%
    \begin{beamercolorbox}[wd=\paperwidth,ht=2.25ex,dp=1ex,right,rightskip=2ex]{footline}%
      \usebeamerfont{page number in head/foot}%
      \insertframenumber{} / \inserttotalframenumber{}
    \end{beamercolorbox}%
  }%
  \vskip0pt%
}


% Style: C:/mytex/2026-Krakow-pres.tex.
% Source: C:/mytex/2026-trade.tex, revised September 19, 2026.
% Focus on the paper's computable countermarket construction.
% Aligned with the moving neutral branch and concise randomness discussion.
% Former financial digressions remain removed.
% All numerical paths are illustrative, not market data.
% 25 frames, including title and section dividers.
% Self-contained; compile twice with pdflatex 2026-ECONVN-pres.tex.
\newcommand{\doiref}[2]{%
  \href{https://doi.org/#1}{\textcolor{blue!65!black}{#2}}}
\newcommand{\sourcefoot}[1]{%
  \par\vfill{\fontsize{6.3}{7.5}\selectfont\color{black!65}Sources: #1\par}}

\title{Computable Countermarkets and the\\Limits of Universal Trading}
\subtitle{A Computable Countermarket for Every Trader}
\author{\texorpdfstring{Karl Svozil\\ITP, TU Wien, Vienna, Austria}%
{Karl Svozil, ITP, TU Wien, Vienna, Austria}}
\date{ECONVN 2026/2027, HUB/AJEB,\\ Ho Chi Minh City, Vietnam\\
12-14 October 2026\\
{\footnotesize \url{https://svozil.github.io/publications/2026-ECONVN-pres.pdf} }
\\
based on
\\
\url{https://doi.org/10.48550/arXiv.2604.13334}
}

\begin{document}

% Frame 1.
\begin{frame}
\titlepage
\end{frame}

% Frame 2.
\begin{frame}{Finite profits do not establish universal profitability}
\small
\begin{block}{A trader can make money over a finite period}
Accurate information, luck, or a continuing trend can produce profits.
Exuberance, innovation, and improvements in production can support a rise.
\end{block}
\begin{exampleblock}{The same kind of bet can also lose}
Information may be false. Exuberance may turn into panic.
Expected gains from innovation or cost cutting may fail to reach shareholders.
\end{exampleblock}
\begin{alertblock}{The claim that fails:}
\par\medskip
{\color{red}\bf ``One trading rule that earns strict profit on every possible
price history.'' (false)}
\par\medskip
Proof by contradiction akin Cantor diagonalization, G\"odel incompleteness, recursive unsolvability of Turing's halting problem: For each total deterministic computable rule, construct a
\emph{countermarket} on which it cannot win---even though prices move at
every date.
\end{alertblock}
\end{frame}


% Frame 3.
\begin{frame}{The trader chooses from the history so far}
\small
\begin{block}{A trading rule}
The rule observes past and current prices and chooses a position before
the next price arrives. It cannot use a future price that it has not observed.
\end{block}
\begin{exampleblock}{Three choices}
\begin{description}
\setlength{\itemsep}{0.55em}
\item[Long:] own shares and gain when their price rises.
\item[Short:] borrow and sell shares, intending to buy them back later;
gain when their price falls.
\item[Wait:] hold neither a long nor a short position; earn zero trading gain.
\end{description}
\end{exampleblock}
\begin{alertblock}{The mathematical setting}
The rule is a total deterministic computable map from finite positive
rational price histories to rational share positions.
It finishes on every input, but no fixed computation deadline is assumed.
\end{alertblock}
\end{frame}

% Frame 4.
\begin{frame}
\hskip 4 cm
\begin{center}
\Large\color{blue}
Countermarket construction
\end{center}
\end{frame}

% Frame 5.
\begin{frame}[shrink=8]{Metamathematical technique: Cantor's diagonal idea}
\small
Diagonal reasoning also appears in the work of G\"odel and Turing.
Suppose there exists a \emph{complete} list containing \emph{all} infinite sequences of zeros and
ones. Its beginning might look like this:
\[
\begin{array}{c|ccccc}
 & \text{bit 1}&\text{bit 2}&\text{bit 3}&\text{bit 4}&\cdots\\\hline
\text{row 1}&\textcolor{red}{0}&1&1&0&\cdots\\
\text{row 2}&1&\textcolor{red}{1}&0&1&\cdots\\
\text{row 3}&0&1&\textcolor{red}{0}&0&\cdots\\
\text{row 4}&1&0&1&\textcolor{red}{1}&\cdots\\
\vdots&\vdots&\vdots&\vdots&\vdots&\ddots
\end{array}
\]
\begin{exampleblock}{Change the diagonal entries}
Read the highlighted bits and change each $0$ to $1$, each $1$ to $0$.
The new sequence begins $1,0,1,0,\ldots$.
\end{exampleblock}
\begin{alertblock}{Why the list cannot be complete?}
Proof by contradiction: Suppose it is complete.
But the new sequence differs from row 1 at bit 1, from row 2 at bit 2,
and from every row at its diagonal position. It is absent from the list.  So, the completeness assumption yields a complete contradiction (relative to \textit{tertium non datur}).
\end{alertblock}
\sourcefoot{Metamathematics; eg, Noson S. Yanofsky,
\doiref{10.2178/bsl/1058448677}{\emph{A Universal Approach to Self-Referential Paradoxes, Incompleteness and Fixed Points} (2003)}.}
\end{frame}

% Frame 6.
\begin{frame}[shrink=5]{Discounted price: value compared with the bank}
\small
\begin{block}{Definition}
A \emph{discounted price} is the asset's price divided by the current
value of one unit of the comparison bank account.
It tells us how many such bank units the asset is worth.
The model treats this account as risk-free.  All money amounts below are in a fixed currency; eg, Vietnamese dong (VND).
\end{block}
\begin{exampleblock}{Example: one bank unit is a deposit of 100 VND}
Suppose the account earns $10\%$ over the year, whereas the share sells at 121 VND.
\par\medskip
\centering
\renewcommand{\arraystretch}{1.5}
\begin{tabular}{@{}lrrc@{}}
 & Bank unit & One share & Discounted price\\\hline
At the start & $100$ & $100$ & $100/100=1$\\
After one year & $110$ & $121$ & $121/110=1.10$
\end{tabular}
\end{exampleblock}
\begin{alertblock}{What the number means}
Selling the share now buys $1.10$ bank units instead of $1$.
The share gained $10\%$ \emph{relative to the bank account}.
Had the share risen only to $110$, its discounted price would still be $1$.
(This adjusts for bank interest, not inflation.)
\end{alertblock}
\end{frame}

% Requires \usepackage{colortbl} in the preamble.

% Frame 7.
\begin{frame}{From Cantor's idea to a countermarket}
\small

\begin{tabular}{@{}
  >{\raggedright\arraybackslash}p{0.485\textwidth}
  @{\hspace{0.03\textwidth}}
  >{\raggedright\arraybackslash}p{0.485\textwidth}
  @{}}
\textcolor{blue}{Cantor's construction}
&
\textcolor{green!50!black}{The trading construction}
\\[0.8em]

\rowcolor{black!15}
Start with a purportedly complete enumeration.
&
Start with a proposed universally profitable rule.
\\[1em]


Use that list itself to construct something that it misses.
&
Use that rule itself to construct a price path on which it fails.
\\[1em]

\rowcolor{black!15}
Changing the selected bit defeats the completeness claim.
&
Opposing each position defeats the universal-profit claim.
\end{tabular}

\medskip
\begin{alertblock}{The common idea is diagonal opposition}
The counterexample depends on the proposed universal object.
Here we construct a market for one fixed trader; we do not need to list
all trading programs.
\end{alertblock}
\end{frame}

% Frame 8.
\begin{frame}{A Computable Countermarket for Every Trader}
\small
\begin{block}{The result, including a moving market}
Fix a positive rational starting price and a rational $0<\epsilon<1$.
For every total deterministic computable trader with rational positions,
there is a fixed computable path with
\[
  P_{t+1}/P_t\in\{1-\epsilon,1+\epsilon\}
\]
on which accumulated discounted gain is nonpositive at every finite date.
\end{block}
\begin{exampleblock}{Every active position loses}
Any nonzero position makes accumulated gain strictly negative.
A trader that always waits earns zero, although prices keep moving.
\end{exampleblock}
\begin{alertblock}{Model and scope}
One discounted asset; self-financing trading through shares and a bank
account; no fees, borrowing limits, or solvency constraints.
This is a worst-case construction, not a model of market equilibrium.
\end{alertblock}
\sourcefoot{KS, base manuscript, revised 19 September 2026,
subsection \enquote{A Computable Countermarket for Every Trader}.}
\end{frame}

% Frame 9.
\begin{frame}{The countermarket: a constructive proof}
\small
\begin{block}{Use the fixed trader's program}
Compute its position $w_t$ from the rational history so far, then set a countermarket price
\[
P_{t+1}=
\begin{cases}
P_t(1-\epsilon),&w_t>0,\\
P_t(1+\epsilon),&w_t<0,\\
P_t(1+\epsilon(-1)^t),\text{ or ``whatever''}&w_t=0.
\end{cases}
\]
At a neutral date (third case), the price rises if $t$ is even and falls if $t$ is odd (arbitrary since $w_t=0$).
\end{block}
\begin{exampleblock}{The market always moves}
Every return has magnitude $\epsilon$. Prices stay positive and rational.
Every nonzero position loses; zero holdings earn zero.
\end{exampleblock}
\begin{alertblock}{Why the construction works}
Adding losses and zeros never gives a positive accumulated gain.
Totality makes each next price computable. For the fixed trader,
any finite part of this fixed path can be generated before trading begins.
\end{alertblock}
\end{frame}

% Frame 10.
\begin{frame}[shrink=5]{Worked example: we construct the adverse prices}
\small
\begin{block}{The trader's rule is fixed first}
The rule is: buy for one period, sell short for the next, then wait.
All prices are discounted prices; gains and losses are measured in
bank-account units.
\end{block}
\begin{exampleblock}{Our construction, step by step}
\begin{enumerate}
\setlength{\itemsep}{0.55em}
\item The trader buys one share for $100$.\newline
We choose $90$ as the next price. Selling the share then loses $10$.
\item The trader borrows a share and sells it for $90$.\newline
We choose $99$ as the next price. Buying it back to return it loses $9$.
\item The trader holds no position.\newline
At date $t=2$, we raise the price by $10\%$ to $108.9$.
Its trading gain is still zero.
\end{enumerate}
\end{exampleblock}
\begin{alertblock}{Why this is a counterexample}
The prices are $100,90,99,108.9$; the gains are $-10,-9,0$.
Total gain is $-19$, although the market moves at every date.
This refutes the rule's claim of strict profit on every permitted history.
\end{alertblock}
\end{frame}

% Frame 11.
\begin{frame}{Why the countermarket is a fixed computable path}
\small
\begin{block}{Put the trader's code inside a price-generating program}
\begin{enumerate}
\setlength{\itemsep}{0.45em}
\item Start with a positive price and the history built so far.
\item Run the trader's program on that history to obtain its position.
\item Choose a small price fall for a long position, a small rise for a
short position, and an alternating-by-date rise or fall for zero holdings.
\item Append that price to the history and repeat.
\end{enumerate}
\end{block}
\begin{exampleblock}{Each requested finite part can be calculated}
The trader returns its position in finite time. Rational positions, expressed
as ratios of integers, can be compared with zero exactly. Each next price is therefore
computable.
\end{exampleblock}
\begin{alertblock}{The market does not have to observe live orders}
For a fixed trader, the generator is fixed. Any finite portion can be
computed before trading begins, although computation may be slow.
\end{alertblock}
\end{frame}

% Frame 12.
\begin{frame}{Every trader has its own countermarket}
\small
\begin{block}{The order of the claim matters}
Proposed guarantee: \enquote{There is one trader that profits on every path.}
\par\medskip
Countermarket result: \enquote{For each proposed trader, construct a path
on which every nonzero position loses.}
\end{block}
\begin{exampleblock}{Another trader may win on that same record}
For one price rise, a trader holding one share gains and a trader shorting
one share loses. A price fall reverses their fortunes.
The construction does not claim to make every other trader lose on that path.
\end{exampleblock}
\begin{alertblock}{The stronger contribution}
A constant discounted price already gives every trader zero trading gain.
Our construction also works when flat paths are excluded: every return
has fixed nonzero magnitude, and every active position of the selected
trader loses. The path need not be typical, an equilibrium, or arbitrage-free.
\end{alertblock}
\end{frame}

% Frame 13.
\begin{frame}{Two sources of randomness}
\small
\begin{block}{Trader side: a private coin chooses positions}
After observing prices, the trader might toss a private coin:
heads means hold one share; tails means short one share.
The program and price history alone no longer determine its next position.
Our deterministic countermarket cannot calculate and oppose that hidden draw.
\end{block}
\begin{exampleblock}{Market side: randomness determines price movements}
The next price may instead be generated randomly, independently of
the trader's private seed. Both sources can occur together.
A passive market does not react to the trader's orders.
\end{exampleblock}
\begin{alertblock}{Randomness alone does not eliminate profits}
Randomizing positions need not destroy an existing edge.
Random market movements can still have positive expected returns.
Zero expected gain requires additional fairness assumptions.
\end{alertblock}
\sourcefoot{KS, subsection \enquote{Worst-Case Value and Private Randomization}.}
\end{frame}

% Frame 14.
\begin{frame}[shrink=8]{An independent fair market gives zero expected gain}
\small
\begin{block}{The market's coin chooses the price movement}
The trader first chooses its position, possibly at random.
An independent fair coin then moves the discounted price from $100$
to $90$ or $110$, each with probability one half.
The market's coin is independent of the past and of the trader's private coin.
\end{block}

\begin{exampleblock}{Average over the market's two equally likely moves}
\centering
\renewcommand{\arraystretch}{1.3}
\setlength{\tabcolsep}{8pt}
\begin{tabular}{@{}lrrc@{}}
Position & Price $90$ & Price $110$ & Expected gain\\\hline
\rowcolor{black!15}
Own one share & $-10$ & $+10$ & $0$\\
Short one share & $+10$ & $-10$ & $0$\\
\rowcolor{black!15}
Hold no position & $0$ & $0$ & $0$
\end{tabular}
\end{exampleblock}

\begin{alertblock}{The result applies to deterministic and randomized traders}
Every position has expected gain zero, however selected.
Repeat independent fair $10\%$ rises or falls from each current price.
Over any fixed finite horizon, expected total gain is zero,
provided gains and losses have finite averages.
\par\smallskip
Individual runs may win or lose. Zero expected gain is not a claim
that realized profits converge to zero in the long run.
\end{alertblock}
\sourcefoot{KS, proposition \enquote{Private Randomization Does Not Create a Guarantee}.}
\end{frame}

% Frame 15.
\begin{frame}
\hskip 4 cm
\begin{center}
\Large\color{blue}Some related ideas
\end{center}
\end{frame}


% Frame 16.
\begin{frame}{Gold: can history reveal the market's rule?}
\small
\begin{block}{Learning in the limit}
After each new observation, a learner proposes a program explaining the
whole sequence. Success means eventually settling on one correct program
and never changing it again. The learner need not know when it has succeeded.
\end{block}
\begin{alertblock}{Gold's negative result}
No algorithm achieves this for \emph{every} computable sequence of zeros
and ones. For each learner, some computable history defeats this requirement.
Failure means no eventual correct identification, rather than a specified
number of losing trades.
\end{alertblock}
\begin{exampleblock}{Restrictions can make learning possible}
For a uniformly effective indexed family of total computable rules,
identification in the limit is possible: keep the first rule consistent
with all observations so far.
What matters is which market rules the learner admits.
\end{exampleblock}
\sourcefoot{E.\,M. Gold,
\doiref{10.2307/2270580}{\emph{Limiting recursion} (1965)};
\doiref{10.1016/S0019-9958(67)91165-5}{\emph{Language identification in the limit} (1967)}.
Related prediction construction: \doiref{10.48550/arXiv.cs/0606070}{Legg, Is there an Elegant Universal Theory of Prediction? (2006)}.}
\end{frame}

% Frame 17.
\begin{frame}[shrink=4]{Turing: will the price ever rise?}
\small

\begin{block}{An artificial market controlled by a program}
Run any program with fixed input, one instruction per day. Start the price at $1$.
\begin{itemize}
\setlength{\itemsep}{0.2em}
\item While the program keeps running, the price stays at $1$.
\item If the program stops, the price becomes $2$ and stays there.
\end{itemize}
\end{block}

\begin{exampleblock}{The price on a specified day is computable}
\centering
\renewcommand{\arraystretch}{1.2}
\begin{tabular}{@{}lccccc@{}}
Program's behaviour & Day 1 & Day 2 & Day 3 & Day 4 & $\cdots$\\\hline
Stops at instruction 3 & $1$ & $1$ & $2$ & $2$ & $\cdots$\\
Never stops            & $1$ & $1$ & $1$ & $1$ & $\cdots$
\end{tabular}
\par\smallskip
\raggedright
For day $100$, simulate $100$ instructions: report $2$ if the program
has stopped, otherwise $1$. This calculation always finishes.
\end{exampleblock}

\begin{alertblock}{But will the price \emph{ever} rise?}
The price rises exactly when the program stops. Deciding this for every
program and input would solve the \emph{Halting Problem}.
No algorithm always finishes with the correct yes-or-no answer.
\end{alertblock}

\sourcefoot{A.\,M. Turing,
\doiref{10.1112/plms/s2-42.1.230}{\emph{On computable numbers} (1936--37)}.
Market construction: the paper.}
\end{frame}

% Frame 18.
\begin{frame}{Rad\'o's Busy Beaver: how long must we wait?}
\small
\begin{block}{A finite maximum with no computable general bound}
Fix a \emph{universal} programming language, capable of simulating any
ordinary program. Include any required input in its program description.
Among programs of at most $n$ bits that eventually stop, take the longest
running time. This is the runtime
\emph{Busy Beaver} value, written $\mathrm{BB}(n)$.
Here $n$ measures program length, not elapsed time.
\end{block}
\begin{exampleblock}{Return to the price $1$ that may become $2$}
A short program can run for an extraordinarily long time before stopping.
Its market can therefore remain flat through a long observation period and
change only later.
\end{exampleblock}
\begin{alertblock}{Why no waiting rule suffices}
No computable deadline based only on program length covers every eventual
halt. Such a deadline would solve Turing's problem: wait until it expires.
Some finite cases can be settled, but they give no general extrapolation rule.
\end{alertblock}
\sourcefoot{T. Rad\'o,
\doiref{10.1002/j.1538-7305.1962.tb00480.x}{\emph{On non-computable functions} (1962)};
G.\,J. Chaitin,
\doiref{10.1007/978-1-4612-4808-8_28}{\emph{Computing the busy beaver function} (1987)}.}
\end{frame}

% Frame 19.
\begin{frame}{From random trials to ML randomness}
\small
\begin{block}{Two levels}
\emph{Process randomness}: how individual bits are generated.
\emph{Martin-L\"of (ML) randomness}: a property of an infinite record
relative to a specified probability law.
\end{block}
\begin{alertblock}{The sufficient condition}
Independent fair bits produce a fair-coin ML-random sequence
\textbf{with probability one}.
\par\smallskip
Independence alone is insufficient if the bits are biased.
Fair one-bit probabilities alone also fail: repeating one fair bit forever
gives a constant record.
\end{alertblock}
\begin{exampleblock}{The practical qualification}
More generally, any computable process law produces records ML-random
for that law almost surely; dependence is allowed.
Finite tests can challenge a source model, but cannot certify the randomness
of its infinite output. Probability one is not a guarantee for every outcome.
\end{exampleblock}
\sourcefoot{P. Martin-L\"of,
\doiref{10.1016/S0019-9958(66)80018-9}{\emph{The definition of random sequences} (1966), Sec.~IV}.}
\end{frame}

% Frame 20.
\begin{frame}{Borel normality is weaker than randomness}
\small
\begin{block}{Correct frequencies at every block length}
A base-$b$ record is \emph{normal} if every word of length $k$ has limiting
frequency $b^{-k}$, for every $k\geq1$.
ML randomness for uniform digits implies normality.
\end{block}
\begin{exampleblock}{Champernowne's decimal}
\[
  0.123456789101112131415\ldots
\]
Concatenate the positive integers. This decimal is normal in base ten,
yet computable: its digits are prescribed by a simple rule.
It is therefore not ML-random for the uniform decimal law.
\end{exampleblock}
\begin{alertblock}{Why verification is difficult}
Normality concerns all block lengths and infinite limits.
No finite frequency check proves it: every finite prefix admits both
normal and nonnormal continuations.
A mathematical construction can establish what sampling alone cannot.
\end{alertblock}
\sourcefoot{D. G. Champernowne,
\doiref{10.1112/jlms/s1-8.4.254}{\emph{The Construction of Decimals Normal in the Scale of Ten} (1933)}.}
\end{frame}

% Frame 21.
\begin{frame}{Ramsey: order appears even without a predictive law}
\small
\begin{block}{An unavoidable pattern}
Join every pair among six objects and color each link red or blue.
Whatever the coloring, three objects have all three links the same color.
This is a basic finite case of Ramsey's theorem.
\end{block}
\begin{exampleblock}{Repeated symbols at equally spaced dates}
Van der Waerden's theorem gives a related result: in a sufficiently long
record with finitely many symbols, some symbol repeats at equally spaced
positions for any prescribed finite number of repetitions.
\end{exampleblock}
\begin{alertblock}{The finite-pattern trap}
The selected dates need not be consecutive, the returns need not be positive,
and the pattern need not continue. Calude and Longo emphasize that large
data sets force some regularities even when the records resist compression.
Such patterns coexist with Martin-L\"of randomness.
\end{alertblock}
\sourcefoot{Graham, Rothschild \& Spencer, \emph{Ramsey Theory} (1990);
\doiref{10.1038/scientificamerican0790-112}{Graham \& Spencer (1990)};
\doiref{10.1007/s10699-016-9489-4}{Calude \& Longo, \emph{The deluge of spurious correlations in big data} (2017)}.}
\end{frame}

% Frame 22.
\begin{frame}{Time reversal as a search for counterexamples}
\small
\begin{block}{Hold the trading rule fixed, change the record}
Test it on the illustrative prices $100,\ 110,\ 120$ and then rerun it
from scratch on $120,\ 110,\ 100$.
At every date it may use only the history available on that run.
Reversing the original orders would give it information from the future.
\end{block}
\begin{block}{If the initial price must stay fixed}
Use $P^{\mathrm{rev}}(t)=P(0)P(T-t)/P(T)$, provided it stays in the class.
This rescales the reversed record; a price-level-sensitive rule may act
differently.
\end{block}
\begin{alertblock}{One failure refutes an all-path claim}
If both records belong to the claimed market class, a nonpositive gain on
either refutes strict profit on every path.
Reversal alone does not prove that a rule must fail on one of the pair.
It supplies a stress test, whereas diagonalization supplies the counterexample.
\end{alertblock}
\sourcefoot{KS, \enquote{Practical Robustness Diagnostics: Time Reversal and Permutations},
Reversal Falsification Test.}
\end{frame}

% Frame 23.
\begin{frame}[shrink=8]{Beyond time reversal: reshuffling the record}
\small
\begin{block}{Algorithmic and random permutations}
A \emph{permutation} rearranges the same finite record of discounted price
levels. Use a fixed algorithm or a random shuffle. Rerun the trader from
scratch, giving it only the past available in the new order.
\end{block}
\begin{alertblock}{A necessary condition for a universal profit guarantee}
A strategy claiming strict profit on \emph{every} permitted path
\textbf{must remain profitable under every reshuffling that stays within
its claimed market class}. The amount of profit may change.
One zero or negative gain refutes that universal claim.
\end{alertblock}
\begin{exampleblock}{What this requirement means}
A strategy exploiting a particular trend may fail when reshuffling destroys
that trend. Such failure limits its scope. Conversely, passing a finite
collection of shuffles cannot establish universal profitability.
\end{exampleblock}
\medskip
\emph{Literary analogy:} Greg Egan's \emph{Permutation City} (1994)
explores the ordering of computational states.
\sourcefoot{Greg Egan, \href{https://www.gregegan.net/PERMUTATION/Permutation.html}{\emph{Permutation City} (1994)};
\href{https://www.gregegan.net/PERMUTATION/FAQ/FAQ.html}{Egan's \emph{Dust Theory FAQ}, questions 1--2}.}
\end{frame}

% Frame 24.
\begin{frame}{Summary of Countermarkets}
\small
\begin{block}{A constructive worst-case result}
For every total deterministic computable rule on rational histories,
there is a positive computable path with a fixed nonzero return magnitude
on which its self-financing gain is nonpositive at every finite horizon.
\end{block}
\begin{alertblock}{What the result does and does not establish}
Every active position loses; waiting earns zero while the market moves.
The path depends on the chosen trader and need not be arbitrage-free.
This is an expository construction, not a new general impossibility theorem.
\end{alertblock}
\begin{exampleblock}{Conditional success remains possible}
Private randomization and random markets do not generally eliminate profits.
Information, luck, or favorable market conditions can produce gains.
Specify the market law, horizon, costs, and benchmark; these gains are
compatible with the failure of an all-path guarantee.
\end{exampleblock}
\end{frame}


% Frame 25.
\begin{frame}
\hskip 4 cm
\begin{center}
\Large\color{blue}Thank you for your attention
\end{center}
\end{frame}

\end{document}
